\begin{abstract}
拟合历史参与总量，并不意味着识别了产生这些总量的行为机制。广泛的信息触达与较低的参与意愿，以及有限的信息触达与较高的参与意愿，都可能解释同一段历史；但在候补、同伴转发和时间冲突共同作用时，两者可能给出不同预测。本文研究这一情境下的预算化过程证据获取。我们保留有限个与历史相容的模拟机制，分别依据过程记录和群体轨迹筛选跨活动共享的阶段参数，再按噪声归一化的最差活动误差分配权重。缓存的组合探针把预期预测方差缩减转化为加权协方差计算，小批量查询通过对假设结果求期望完成条件化，无需重新运行模拟器。在 V1 可审计的 60 候选开发搜索中，归一化轨迹误差 AUC 下降 \DevReduction\%，冻结合成测试 AUC 为 \TestAUC；GPU 配对微基准显示，协方差评分比所提供的显式枚举实现快 \FrozenSpeed\ 倍。然而，24 个独立合成世界的评测揭示了一个重要限制：所选校准分布的有效机制数只有 \ESS，同先验取证、匹配单点估计以及两项取证消融的表现几乎一致，未知规则诊断也会漏检较弱偏移。这些结果给出了高效、可复现的实现，同时表明：群体预测的改善，可能压缩主动取证所依赖的机制分歧。另行开展的 30 配置扩展保留了非坍缩先验，在构造交互压力测试上恢复了同先验取证增益，但与参数信息增益持平。本文仅提供合成证据，不主张真实数据成效、官方榜单排名或一般性的次模保证。
\end{abstract}
